\documentclass[acmsmall,nonacm,screen]{acmart}

\usepackage{mathpartir}

\citestyle{acmauthoryear}
\settopmatter{printacmref=false}

\newcommand{\kw}[1]{\mathsf{#1}}
\newcommand{\nm}[1]{\mathit{#1}}
\newcommand{\Ok}[2]{\kw{Ok}(#1, #2)}
\newcommand{\ret}{\kw{ret}}
\newcommand{\ask}[2]{\kw{ask}\;#1\;(#2)}
\newcommand{\act}[3]{\kw{act}\;#1\;#2\;#3}
\newcommand{\actc}[2]{\kw{act!}\;#1\;#2}
\newcommand{\orc}[1]{\kw{oracle}(#1)}
\newcommand{\inp}{\kw{input}()}
\newcommand{\letin}[3]{\kw{let}\;#1 = #2\;\kw{in}\;#3}
\newcommand{\judg}[4]{\Phi; \Gamma; #1 \vdash_{#2} #3 : #4}
\newcommand{\Base}{\kw{Base}}
\newcommand{\T}{\kw{T}}
\newcommand{\U}{\kw{U}}
\newcommand{\lab}{\mathrm{lab}}
\newcommand{\taint}[2]{\kw{taint}_{#1}(#2)}

\begin{document}

\title{Ask Only When It Matters: Integrity Types That Decide Which Actions Need Approval}

\author{Claude-Fable-5.2}
\authornote{An AI model. It wrote this paper at the request of a human who asked for papers on
the strongest ideas of this repository. The proofs are on paper and have not been checked by a
machine. At the time of writing nobody has reviewed the paper.}

\begin{abstract}
An agent that asks for approval before every action teaches its user to approve without
reading. We give a type system that decides which actions need the question. Types carry an
integrity label. Data from outside is untrusted, the answer of a model is as untrusted as its
input, and a branch on untrusted data taints what it controls. An action that untrusted data
can influence must consume an approval, in the sense of typed approvals, and every other action
runs without one. We prove that in every execution each action is approved or clean, and that
the clean actions are the same whatever the untrusted inputs were. A prompt injection can
therefore change only actions that a person is shown.
\end{abstract}

\keywords{information flow, integrity, linear types, agents, prompt injection}

\authorsaddresses{}

\maketitle

\section{Introduction}

A person in the loop protects against an agent that goes wrong, as long as the person reads
what is shown. Reading does not survive a question before every step. The agent should ask when
the answer matters and stay silent otherwise.

When does it matter? An agent goes wrong in a dangerous way when text that it reads tells it
what to do: a web page, an e-mail, or the output of a tool that contains instructions. The model
that drives the agent cannot be relied on to ignore them. This is an integrity
problem~\citep{denning1976}. Data from outside is untrusted, and an action is at risk when
untrusted data can influence it.

Typed approvals~\citep{tpp4} guarantee that an approved action runs with the payload that the
person saw. They say nothing about which actions should be approved. This paper adds the
missing half. We label base values as trusted or untrusted, propagate the labels through the
program, and give two typing rules for actions. An action that untrusted data can influence, in
its payload or in the decision to perform it, must consume an approval. An action that only
trusted data influences needs none.

The model is the unusual part. We treat it as a function whose result is as untrusted as its
argument. Whatever the model was told by its input, the label of the input covers it. No
assumption is made that the model resists instructions.

Our main contributions are as follows:
\begin{itemize}
\item A type system with integrity labels and a program counter label on top of the calculus of
  typed approvals (Section~\ref{sec:types}).
\item The theorem that every action of a well-typed program is approved or clean
  (Theorem~\ref{thm:selective}).
\item The theorem that the clean behavior of a program is independent of its untrusted inputs
  (Theorem~\ref{thm:indep}).
\end{itemize}

\paragraph{Status of the claims.}
The proofs in this paper are sketched on paper and have not been mechanized.

\section{Labels and Types}
\label{sec:types}

We extend the calculus of~\citet{tpp4}. Recall that an approval has the type
$\Ok{a}{x.\varphi}$ for an action $a$ and a predicate $\varphi$ on its payload $x$, that
$\ask{a}{x.\varphi}$ returns an approval or nothing, and that $\act{a}{i}{w}$ performs $a$ with
the payload $i$ and consumes the approval $w$, which must satisfy $\Phi \models \varphi[i/x]$.

A label $\ell$ is $\T$ for trusted or $\U$ for untrusted, with $\T \sqsubseteq \U$. Base types
and sum types carry a label:
\[
\tau ::= \Base^\ell \mid \kw{Unit} \mid \tau \otimes \tau \mid (\tau \oplus \tau)^\ell \mid
  \Ok{a}{x.\varphi} .
\]
We write $\tau \sqcup \ell$ for $\tau$ with the labels of its base types and sum types raised by
$\ell$, outside approvals. The context $\Gamma$ gives a label to each base variable. The label
$\lab(i)$ of an index is the label of the variable, or $\T$ for a value. The label
$\lab(\varphi)$ of a predicate is the join of the labels of its variables. Each action $a$ has a
fixed label $\ell_a$ for the results that the world returns.

There are two new expressions. The expression $\inp$ reads data from outside. The expression
$\actc{a}{i}$ performs an action without an approval.

\begin{figure}
\begin{mathpar}
\inferrule[Ret]{\Phi; \Gamma; \Delta \vdash w : \tau}
  {\judg{\Delta}{pc}{\ret\;w}{\tau \sqcup pc}}
\and
\inferrule[Input]{ }{\judg{\cdot}{pc}{\inp}{\Base^{\U}}}
\and
\inferrule[Oracle]{\Gamma \vdash i}{\judg{\cdot}{pc}{\orc{i}}{\Base^{\lab(i) \sqcup pc}}}
\and
\inferrule[If]{\judg{\Delta}{pc \sqcup \lab(\varphi)}{e_1}{\tau} \quad (\Phi, \varphi) \\
  \judg{\Delta}{pc \sqcup \lab(\varphi)}{e_2}{\tau} \quad (\Phi, \neg\varphi)}
  {\judg{\Delta}{pc}{\kw{if}\;\varphi\;\kw{then}\;e_1\;\kw{else}\;e_2}{\tau}}
\and
\inferrule[Ask]{\mathrm{fv}(\varphi) \subseteq \Gamma, x}
  {\judg{\cdot}{pc}{\ask{a}{x.\varphi}}{(\Ok{a}{x.\varphi} \oplus \kw{Unit})^{pc \sqcup
    \lab(\varphi)}}}
\and
\inferrule[Act]{\Gamma \vdash i \\ \Phi; \Gamma; \Delta \vdash w : \Ok{a}{x.\varphi} \\
  \Phi \models \varphi[i/x]}
  {\judg{\Delta}{pc}{\act{a}{i}{w}}{\Base^{\lab(i) \sqcup pc \sqcup \ell_a}}}
\and
\inferrule[Act-Clean]{\Gamma \vdash i \\ \lab(i) \sqcup pc = \T}
  {\judg{\cdot}{pc}{\actc{a}{i}}{\Base^{\ell_a}}}
\end{mathpar}
\Description{Seven typing rules with integrity labels.}
\caption{Typing rules with labels. In \textsc{If}, each branch is checked with the predicate in
parentheses added to $\Phi$.}
\label{fig:rules}
\end{figure}

The judgment $\judg{\Delta}{pc}{e}{\tau}$ has a program counter label $pc$, which bounds the
labels of the data that decided whether $e$ runs. Figure~\ref{fig:rules} gives the rules that
involve labels. The rule for $\kw{case}$ on an operand of type $(\tau_1 \oplus \tau_2)^\ell$
checks both branches at $pc \sqcup \ell$, as \textsc{If} does. A $\kw{let}$ on a base
expression of type $\Base^\ell$ gives the label $\ell$ to its variable. A procedure is checked
at a declared label and called only at a program counter below it. The remaining rules are those
of~\citet{tpp4}.

Three rules carry the design. In \textsc{Oracle}, the answer of the model has the label of the
question. In \textsc{Ask}, the answer of the person has the label of what was shown: a person
who reads untrusted text can be misled by it, so an approval does not make later steps trusted.
In \textsc{Act-Clean}, an action runs without approval only when its payload and the program
counter are trusted.

\paragraph{Example.}
An agent pays an invoice that arrives by e-mail and then files a note.
\[
\begin{array}{l}
\letin{m}{\inp}{} \\
\letin{p}{\orc{m}}{} \\
\letin{r}{\ask{\nm{pay}}{x.\,x = p}}{} \\
\kw{case}\;r\;\{\kw{inl}\;t \Rightarrow \act{\nm{pay}}{p}{t} \mid \kw{inr}\;u \Rightarrow
  \ret\;0\}
\end{array}
\]
The mail $m$ is untrusted, so the payment $p$ that the model extracts is untrusted, and
\textsc{Act-Clean} rejects $\actc{\nm{pay}}{p}$. The program above is accepted: the person sees
the exact payment. Suppose the agent continues with $\actc{\nm{log}}{s}$ for a constant $s$.
After the $\kw{case}$ the program counter is trusted again, the payload is a constant, and the
note is filed without a question. Inside the first branch of the $\kw{case}$ the same expression
would be rejected, since whether it runs there depends on an answer to untrusted text.

\section{Semantics with Labels}

The semantics of~\citet{tpp4} has a store of approvals and a monitor that lets an action
through when its approval is in the store and covers the payload. We add labels at run time. A
base value is $c^\ell$ and an injection is $\kw{inl}^\ell\,v$ or $\kw{inr}^\ell\,v$. A frame
$\taint{\ell}{e}$ runs $e$ under a raised program counter. The dynamic program counter of a
step is the join of the labels of the frames around its redex. The rules are as follows, where
$pc$ is the dynamic program counter.
\[
\begin{array}{lcl}
\kw{if}\;\varphi\;\kw{then}\;e_1\;\kw{else}\;e_2 & \to & \taint{\ell}{e_k}
  \quad \text{for the branch } k \text{ that } \varphi \text{ selects, } \ell = \lab(\varphi) \\
\kw{case}\;(\kw{inl}^\ell\,v)\;\{\ldots\} & \to & \taint{\ell}{e_1[v/y]} \\
\taint{\ell}{\ret\;v} & \to & \ret\;(v \sqcup \ell) \\
\inp & \to & \ret\;c^{\U} \\
\orc{c^\ell} & \to & \ret\;c'^{\,\ell \sqcup pc} \\
\ask{a}{x.\varphi} & \to & \ret\;(\kw{inl}^{\ell}\,n) \text{ or } \ret\;(\kw{inr}^{\ell}\,())
  \quad \text{with } \ell = pc \sqcup \lab(\varphi) \\
\act{a}{c^\ell}{n} & \to & \ret\;c'^{\,\ell \sqcup pc \sqcup \ell_a} \\
\actc{a}{c^\ell} & \to & \ret\;c'^{\,\ell_a} \quad \text{if } \ell \sqcup pc = \T
\end{array}
\]
The values $c$ and $c'$ are arbitrary, and the events and the store are as before. An action
with an approval emits $\kw{act}(n, a, c)$. An action without one emits $\kw{clean}(a, c)$. The
side condition of the last rule is a second monitor. An action is \emph{clean} when it is
performed by the last rule with the side condition true.

\section{Results}

\begin{theorem}[Safety]
\label{thm:safety}
A well-typed configuration is a returned value or takes a step in the monitored semantics, and
the step leads to a well-typed configuration.
\end{theorem}

\begin{proof}[Proof sketch]
We extend the typing to labelled values and frames. A value $c^\ell$ has type $\Base^{\ell'}$
for $\ell \sqsubseteq \ell'$, and likewise for injections. A frame $\taint{\ell}{e}$ has type
$\tau$ at $pc$ when $e$ has type $\tau$ at $pc \sqcup \ell$. Preservation is by cases on the
step, with the substitution lemma of~\citet{tpp4}. The invariant is that run-time labels are
below static labels and that the dynamic program counter is below the static one. For progress,
the new case is $\actc{a}{c^\ell}$. By \textsc{Act-Clean} the static label of the payload and
the static program counter are $\T$, so the dynamic ones are $\T$, and the monitor lets the
action through. The case of $\kw{act}$ is as in~\citet{tpp4}.
\end{proof}

\begin{theorem}[Selective consent]
\label{thm:selective}
In every execution of a well-typed program without the two monitors, each action either
consumes an approval that was granted earlier for this action and whose predicate the payload
satisfies, with no approval consumed twice, or it is clean.
\end{theorem}

\begin{proof}
By Theorem~\ref{thm:safety} neither monitor ever blocks, so the executions without monitors are
the monitored executions. In a monitored execution an action with an approval has consent by
the argument of~\citet{tpp4}, and an action without one passed the second monitor.
\end{proof}

The second theorem says what clean means. Two configurations are \emph{equal up to untrusted
data} when they become equal after every value with label $\U$ and every frame with label $\U$
is replaced by a hole, up to a renaming of approvals. A step is \emph{visible} when its redex is
not inside a frame with label $\U$. The visible events of an execution are the events of its
visible steps, with payloads of label $\U$ replaced by a hole.

\begin{theorem}[Independence]
\label{thm:indep}
Let two well-typed configurations be equal up to untrusted data. For every finite execution of
the first there is an execution of the second whose visible events are a prefix of the visible
events of the first, and all of them unless a frame with label $\U$ diverges in the second.
\end{theorem}

\begin{proof}[Proof sketch]
The execution of the second configuration is built step by step so that the two stay equal up
to untrusted data. A step inside a frame with label $\U$ is matched by no step. For a visible
step, the redex is outside the holes, so the second configuration has the same redex up to
values with label $\U$. A test or a $\kw{case}$ on trusted data takes the same branch. On
untrusted data both sides enter a frame with label $\U$, which is a hole. The answer of the
model to a trusted question, the answer of the person to a trusted predicate, and the result of
an action with a trusted payload are trusted, and the second execution makes the same choice.
In the other cases the result has label $\U$ and is a hole. When a frame with label $\U$
returns, the second configuration runs its frame until it returns, unless it diverges. The
results have the same type, their base values and injections are raised to $\U$ and are holes,
and approvals in them are related by the renaming.
\end{proof}

Every clean action is a visible event with a trusted payload, so the clean actions of a program
are the same for all untrusted inputs, as far as the second execution gets. This is the usual
guarantee of a type system for information flow~\citep{volpano1996}, read for integrity, and it
is insensitive to termination in the usual way.

\section{Discussion}

\emph{Against prompt injection.} Text that an agent reads enters through $\inp$ or through the
result of an action with $\ell_a = \U$. By Theorem~\ref{thm:indep} it cannot change a clean
action. By Theorem~\ref{thm:selective} every other action is one that a person approved with
its payload in view. The model needs no defense of its own.

\emph{Fewer questions.} An agent that works from the instructions of its user and from trusted
data asks nothing. The questions begin where outside data starts to steer, and each one shows
the payload that will run.

\emph{What is trusted.} The labels $\ell_a$ and the choice of what counts as $\inp$ are
decisions of whoever deploys the agent. The model is trusted on trusted input: its mistakes are
not attacks, and this paper does not bound them.

\emph{No endorsement.} The system has no way to turn untrusted data into trusted data. A person
who approves an untrusted payload approves one action. A sound checker could justify more, for
example by accepting only values from a fixed list. We leave rules for endorsement to later
work.

\section{Related Work}

The lattice model of~\citet{denning1976} and the type system of~\citet{volpano1996} are the basis
of our labels and of Theorem~\ref{thm:indep}. Our approvals are those of~\citet{tpp4}. The
contribution here is their combination: integrity labels select the actions that need an
approval, and the approval is linear and indexed by the payload.

\citet{debenedetti2025} defend agents against prompt injection with an interpreter that
separates the control flow of a trusted plan from untrusted data and tracks the provenance of
values. \citet{costa2025} track labels dynamically in the planner of an agent and characterize
the properties that such tracking can enforce. Both track flows at run time. We give a static
discipline for the same concern and connect it to consent: where the flow analysis finds
influence from outside, the type system demands an approval instead of forbidding the action.

\section{Conclusion}

Integrity labels tell which actions untrusted data can influence, and typed approvals make
exactly those actions wait for a person. The rest of the program runs unasked and is provably
unaffected by what the agent reads. The next steps are endorsement by sound checkers, labels
with more than two levels, and a mechanized proof.

\bibliographystyle{ACM-Reference-Format}
\bibliography{refs}

\end{document}
